% NichidaiMasterExam_R8_2nd_2
\begin{tcolorbox} %%R8_2nd_2
	$\fbox{2}$ $n$次単位行列を$E_n$、行列$X$の転置を$\tr{X}$とかく。次の問いに答えよ。
\begin{enumerate}
	\item ベクトル$\bm{v}= \frac{1}{\sqrt{5}}\begin{pmatrix}-1\\ 1\end{pmatrix}$に対して行列$A:= E_2- 2\bm{v}\tr{\bm{v}}$を定める。
\begin{enumerate_alpha}
	\item 行列$A$を求めなさい。
	\item 行列$A$の固有値をすべて求めなさい。
\end{enumerate_alpha}
	長さが1である実$n$ベクトル$\bm{w}\in \R^n$に対して$B:= E_n- 2\bm{w}\tr{\bm{w}}$を定める。
\begin{enumerate_alpha}
	\item $\tr{\bm{w}}\bm{w}= 1$が成り立つことを示せ。
	\item $\bm{w}$は$B$の固有ベクトルであることを示しなさい。また対応する固有値を答えなさい。
	\item 直交行列の定義を述べなさい。また$B$は直交行列であることを示しなさい。
\end{enumerate_alpha}
\end{enumerate}
\end{tcolorbox}

\begin{souanswer} %R8_2nd_2_ans
\begin{enumerate}
	\item 
\end{enumerate}
\end{souanswer}